\section{Setup}

\subsection{Model}

We work in the \textbf{known-i.i.d.} model of online bipartite matching.
There is a set \(R\) of \(n\) \textbf{offline resources} and a
\textbf{type graph} on \(r\) \textbf{online types}, where type \(\ell\)
has a neighborhood \(N(\ell)\subseteq R\) (the resources able to serve
it). An \textbf{instance} is generated by drawing \(m\) arrivals i.i.d.
from a distribution \(p\) over the \(r\) types; the \(i\)-th arrival
reveals its type (hence its neighborhood) and must be matched,
immediately and irrevocably, to a currently unmatched resource in its
neighborhood, or left unmatched. The type graph and \(p\) are known to
the algorithm; the realized arrival sequence is not. The classical
setting is \(r=n\) and \(m=n\) (each offline vertex a distinct type); we
also use \emph{few-types} instances (\(r\ll n\)), where a prefix
distribution-test is statistically meaningful.

We measure the \textbf{competitive ratio}
\(\rho(\mathrm{ALG}) = \mathbb{E}[|\mathrm{ALG}|] /
\mathbb{E}[|\mathrm{OPT}|]\), where \(\mathrm{OPT}\) is a maximum
matching of the \emph{realized} instance, computed exactly by
Hopcroft--Karp \cite{hopcroftkarp1973}. All reported ratios are
Monte-Carlo estimates with 95\% confidence intervals (Section 2.5). The
advice-free baseline we compare against throughout is KVV Ranking
\cite{kvv1990ranking}, whose ratio we denote \(\rho_{\mathrm{base}}\)
and call the \textbf{baseline strength}.

\subsection{Algorithms}

All the greedy-type algorithms share one primitive. Given a rank (a
total order) on the resources, \textbf{GreedyWithPermutation} matches
each arrival to its available neighbor of \emph{minimum rank}. The
advice-free baselines are three instances of this primitive together
with the two stochastic-matching algorithms:

\begin{itemize}
\tightlist
\item
  \textbf{SimpleGreedy} --- GreedyWithPermutation under the identity
  rank.
\item
  \textbf{Ranking} \cite{kvv1990ranking} --- GreedyWithPermutation under
  a uniformly random rank; the baseline.
\item
  \textbf{Feldman} \cite{feldman2009online} and \textbf{Jaillet--Lu}
  \cite{jailletlu2014online} --- the known-i.i.d. stochastic-matching
  algorithms (a max-flow blue/red decomposition, and a cap-3 LP realized
  by max-flow with per-type list distributions, respectively).
\end{itemize}

\textbf{Degree predictions (object A).} MinPredictedDegree (MPD)
\cite{aci2022mpd} is GreedyWithPermutation ranked by ascending
\emph{predicted degree} \(\mu\in\mathbb{R}^{R}\): it matches to the
available neighbor predicted to be rarest. With \(\mu\) constant it
reduces to Ranking; with \(\mu\) the true realized degrees it is the
consistency ceiling (MinDegree). Following ACI, we also form the
\textbf{augmentations} Feldman(MPD) and JailletLu(MPD), which use the
MPD rank only as the tie-break of the greedy fallback, so the
worst-case-optimal base matching carries the load.

\textbf{Type-histogram advice (object B) and test-and-fallback.} The
advice is a count vector \(\hat c\in\mathbb{Z}_{\ge 0}^{r}\) over types
(equivalently a distribution \(q=\hat c/\|\hat
c\|_1\)). From \(\hat c\) we build an \emph{advice graph}
(\(\hat c_\ell\) copies of type \(\ell\)), compute a maximum matching
\(\hat M\) on it, and record, per type, the resources \(\hat M\) uses.
\textbf{FollowPrediction} (blind Mimic) routes every arrival according
to \(\hat M\). \textbf{TestAndMatch}
\cite{choo2024imperfect,bem2026testmatch} instead Mimics a sublinear
prefix of \(k\) arrivals, estimates the \(\ell_1\) distance between the
prefix type distribution and \(q\), and --- if the estimate is at most a
threshold \(\tau\) --- continues to Mimic, else falls back to Ranking.
The two variants differ only in an early-exit rule and in \(\tau\)
(Choo: \(\tau =
2(\hat n/n - \beta)\); BEM: \(\tau = 2(\hat n/n)(1-\beta)/(1+\beta)\)),
where \(\beta\) is the worst-case baseline ratio. Faithfulness note: the
papers' \(\ell_1\) tester \cite{jiao2018l1} has no deployable
implementation and the authors themselves fall back to an
empirical-\(\ell_1\) proof-of-concept; we do the same, and mark it as
such.

\textbf{The combiner (a benchmarked baseline, not a contribution).} We
port the blind-follow-with-switching combiner of Chłędowski et
al.~\cite{chledowski2021caching} --- run FollowPrediction and Ranking as
shadow experts and follow the leader with hysteresis --- to
contextualize test-and-fallback. We benchmark it; we do not claim it.

\subsection{Prediction objects and the structured error
model}

To compare algorithms that consume \emph{different} prediction objects
(degree vectors \(\mu\) vs type histograms \(\hat c\)), we corrupt each
object with a per-family knob and report the families in parallel; this
heterogeneity is intrinsic and is precisely why no unified benchmark
existed.

For \textbf{degree predictions} we perturb the true degrees
\(\mu^\star\) with four structured error models, each returning both an
\(\ell_1\) \emph{magnitude} and a normalized \textbf{Kendall-\(\tau\)
order error} in \([0,1]\): \emph{random-flip} (replace a fraction of
entries by uniform noise), \emph{systematic-bias} (a monotone rescale
--- order-preserving, hence Kendall-\(\tau\equiv 0\) by construction),
\emph{adversarial} (reflect a fraction of entries, the most
order-damaging), and \emph{distribution-drift} (blend toward degrees
from an independently drawn graph of the same family). Because MPD
depends on \(\mu\) \emph{only through the induced order}, the
order/magnitude distinction is not cosmetic --- it is the axis Section 4
is about.

For \textbf{type-histogram advice} we perturb the realized counts
\(c^\star\) toward a concentrated (spiky) random target by a fraction
\(\eta\in[0,1]\) (\(\eta=0\): perfect advice, \(\hat c=c^\star\);
\(\eta=1\): advice unrelated to the truth), reporting the induced
\(\ell_1(p,q)\) as the advice-error axis for the test-and-fallback
experiments.

\subsection{Consistency and
robustness}

For a prediction-consuming algorithm we call \(\rho\) under
\emph{perfect} advice its \textbf{consistency} and its worst-case
\(\rho\) over \emph{adversarial} advice its \textbf{robustness}. A
prediction algorithm is useful only if it is consistent \emph{and} never
falls (much) below \(\rho_{\mathrm{base}}\); the tension between the
two, and its fundamental limits for the test-and-fallback class, are the
subject of Sections 5 and 7.

\subsection{Experimental methodology and
reproducibility}

The harness is a small dependency-light Python codebase
(NumPy/SciPy/NetworkX). We draw all randomness from independent,
reproducible streams via NumPy's spawn-based generators, and use
\textbf{paired trials}: within a comparison, every algorithm and every
error level reuses the same graph, arrival sequence, \(\mathrm{OPT}\),
and tie-break seed, so differences are attributable to the prediction
alone. Confidence intervals are 95\% normal-approximation half-widths
over trials. The Phase-2 portion reproduces a subset of Borodin,
Karavasilis \& Pankratov \cite{borodin2018experimental}; every figure
and table in this paper is regenerated from a fixed seed by a single
script, listed with its output in Appendix A. Most of those scripts are
self-contained and run on a clean checkout; the real-graph and
real-trace results of Sections 6 and 8 first require external data,
whose provenance and retrieval are recorded in Appendix A.4.

One implementation detail is load-bearing for that claim. Three
preprocessing steps --- Feldman's and Jaillet--Lu's suggested matchings,
and the serving advice \(b\)-matching --- take a maximum flow and
consume its \emph{decomposition}, which, unlike the flow value, is not
unique; the solver picks one by iterating node containers. We therefore
label those networks' nodes with integers rather than strings, since
Python randomises string hashes per process and string labels made the
chosen decomposition, and every number derived from it, differ between
runs of the same script.
